\documentclass{umk-matematika}
\begin{document}
\maketitlepage
    {Практическая работа 13}
    {}
    {}

\textbf{Задача 1.} Сформулируйте прямую теорему о трёх перпендикулярах.\par\vspace{8pt}
\textbf{Задача 2.} Сформулируйте обратную теорему о трёх перпендикулярах.\par\vspace{8pt}
\textbf{Задача 3.} Из точки A к плоскости α проведены перпендикуляр AB и наклонная AC. Что является проекцией наклонной AC на плоскость α?\par\vspace{8pt}
\textbf{Задача 4.} В треугольнике ABC угол C прямой. Отрезок CD — перпендикуляр к плоскости ABC. Назовите проекцию наклонной AD на плоскость ABC.\par\vspace{8pt}
\textbf{Задача 5.} Из точки M к плоскости α проведён перпендикуляр MO = 8 см и наклонная MA = 17 см. Проекция наклонной OA = 15 см. Перпендикулярна ли прямая a, проведённая в плоскости через точку A к OA, наклонной MA?\par\vspace{8pt}
\textbf{Задача 6.} В кубе $ABCDA\_1B\_1C\_1D\_1$ докажите, что прямая BD перпендикулярна прямой AC₁.\par\vspace{8pt}
\textbf{Задача 7.} Из точки A к плоскости проведён перпендикуляр AB = 9 см и наклонная AC = 15 см. В плоскости через точку C проведена прямая, перпендикулярная BC. Найдите расстояние от A до этой прямой.\par\vspace{8pt}
\textbf{Задача 8.} Из вершины A квадрата ABCD проведён перпендикуляр AK к плоскости квадрата. Сторона квадрата = 6 см, AK = 8 см. Найдите расстояние от точки K до прямой BD.\par\vspace{8pt}
\textbf{Задача 9.} Из точки M к плоскости α проведены перпендикуляр MO и две наклонные MA и MB. Известно, что MA = 10 см, MB = 17 см, а их проекции относятся как 6:15. Найдите расстояние от M до плоскости.\par\vspace{8pt}
\textbf{Задача 10.} Строительная мачта высотой 12 м укреплена тросом, закреплённым на земле на расстоянии 5 м от основания мачты. Перпендикулярна ли линия крепления троса к земле самой мачте, если трос натянут под прямым углом к линии на земле, соединяющей основание мачты с точкой крепления?\par\vspace{8pt}
\newpage
\section*{Ответы}

\begin{enumerate}
  \item Ответ: прямая теорема.
  \item Ответ: обратная теорема.
  \item Ответ: отрезок BC.
  \item Ответ: отрезок AC.
  \item Ответ: да, перпендикулярна.
  \item Ответ: BD ⟂ AC₁.
  \item Ответ: 15 см.
  \item Ответ: $\sqrt{82}$ см.
  \item Ответ: $MO = 8$ см.
  \item Ответ: да, перпендикулярна (мачта ⟂ земле, линия ⊂ земле).
\end{enumerate}
\end{document}
